
\def\PUDcodigo{ECA201}

\if\COMPILAR0
   \def\PUDcodacad{04507.52}
\else
   \def\PUDcodacad{04506.58}
\fi

\def\PUDdisciplina{Manufatura Integrada por Computador}

\def\PUDsemestre{Opt}

\def\PUDhoras{80}

%NOVOS ---------------------------------------------------
\def\PUDhorasTeoria{40}
\def\PUDhorasPratica{40}

\def\PUDinterECA{---}

\def\PUDatuaisECA{---}

\def\PUDrecursos{Projetor multimídia; Quadro branco e pincel.}

%----------------------------------------------------------

\def\PUDcreditos{4}

\def\PUDprerequisito{---}

\def\PUDpreacad{---}

\def\PUDobjetivos{Reconhecer as máquinas com Comando Numérico Computadorizado.  Conhecer a linguagem de máquinas NC.  Conhecer um sistema CAD/CAM: suas vantagens e aplicações.  Identificar uma célula de manufatura flexível.  Reconhecer um sistema integrado de manufatura por computador e avaliar suas vantagens e suas desvantagens.}

\def\PUDementa{
%Programação CN;  Sistema CAD/CAM;  Descrição do sistema CAD/CAM;  Software de CAD/CAM - MasterCam;  Comandos para geração de primitivas geométricas;  Comandos para a edição de um desenho;  Projetar através do CAD;  Desenho de ferramentas;  Desenho da peça a ser usinada;  Gerar e transmitir o programa NC para a máquina;  Usinagem;  
Definição e histórico do CIM;  Célula de manufatura flexível (FMS);  Componentes CIM, integração de dados e operações;  Gerenciamento da informação dos componentes CIM;  Procedimentos e gerenciamento de projeto para desenvolver uma estratégia CIM;  Definição das cadeias de processo CIM;  Software de aplicações (ERP, MES);  Casos CIM.}

\def\PUDconteudo{ %& Programação CNC. Reconhecer o torno Comando Numérico Computadorizado. Elaborar programas aplicados a torno CNC e fresadora CNC.  Analisar o funcionamento do torno CNC. Executar operações fundamentais na usinagem de peças no torno CNC. 
%\\ 
 %& Sistema CAD/CAM.      Descrição do sistema CAD/CAM.      Software de Cad/Cam -  MasterCam.      Comandos para geração de primitivas geométricas.      Comandos para a edição de um desenho.      Projetar através do CAD.      Desenho de  ferramentas.      Desenho da peça a ser usinada.      Gerar o programa NC.      Transmissão do programa gerado para o trono CNC.      Usinagem da peça. 
%\\ 
 &  Introdução ao CIM.      Conceitos.      Histórico.      Sistemas Produtivos de Manufatura.      PCP informatizado. 
\\ 
 & Tecnologia CIM.      Elementos do CIM.      Modelo Y.      Tecnologias de Implementação.      ERP (Planejamento de Recursos Empresariais).      FMS (Sistemas Flexíveis de Manufatura).      Noções de Robótica. 
\\ 
 & Prática em CIM.       Planta CIM: Características e Aplicações.      Robótica Aplicada (FMS): - Visão Artificial;   Robô FANUC;  CNC Romi. }

\def\PUDmetodo{Aulas expositórias que podem ser teóricas e/ou práticas, onde as práticas no laboratório serão marcadas no decorrer da disciplina.}

\def\PUDavalia{A avaliação será feita com aplicação de provas teóricas e/ou práticas, além da possibilidade de inclusão de trabalhos e seminários no decorrer da disciplina.}

\def\PUDbibbasica{ &  \href{www.biblioteca.ifce.edu.br/index.asp?codigo_sophia=62265}{Fitzpatrick}, Michael.  Introdução à usinagem com CNC : comando numérico computadorizado.  Porto Alegre - RS.  2013.  365 p.  ISBN: 9788580552515.
\\ 
 &  \href{www.biblioteca.ifce.edu.br/index.asp?codigo_sophia=62087}{Affonso}, Luiz Otávio Amaral.  Equipamentos mecânicos: análise de falhas e solução de problemas.  3º ed.  Rio de Janeiro: Qualitymark, 2014.  387p.  ISBN: 9788541400367. 
\\ 
 &   \href{www.biblioteca.ifce.edu.br/index.asp?codigo_sophia=62510}{Souza}, Adriano Fagali de.  Engenharia integrada por computador e sistemas CAD/CAM/CNC: princípios e aplicações.  2º ed.  São Paulo - SP.  Artliber.  2013.  358 p.  ISBN: 9788588098909. }

\def\PUDbibcomplementar{ &  \href{www.biblioteca.ifce.edu.br/index.asp?codigo_sophia=62912}{Groover}, Mikell P.  Automação Industrial e Sistemas de Manufatura.  3º ed.  São Paulo - SP.  Pearson, 2015.  581 p.  ISBN: 9788576058717.
\\ 
 &  \href{www.biblioteca.ifce.edu.br/index.asp?codigo_sophia=28300}{Rosário}, João Maurício.  Princípios de mecatrônica.  São Paulo (SP): Pearson Prentice Hall, 2005.  356p.  ISBN: 9788576050100.
\\ 
 &  \href{www.biblioteca.ifce.edu.br/index.asp?codigo_sophia=24218}{Affonso}, Luiz Otávio Amaral.  Equipamentos mecânicos: análise de falhas e solução de problemas.  2º ed.  Rio de Janeiro: Qualitymark, 2006.  321p.  ISBN: 8573036346.
\\
 &  \href{www.biblioteca.ifce.edu.br/index.asp?codigo_sophia=57490}{Paranhos Filho}, Moacyr.  Gestão da Produção Industrial.  E-book.  Intersaberes.  346p.  ISBN: 9788565704847.
\\
 &  \href{www.biblioteca.ifce.edu.br/index.asp?codigo_sophia=65223}{Robert} E. Bateman et al.  Simulação de sistemas: aprimorando processo de logística, serviços e manufatura.  Rio de Janeiro, RJ: Elsevier, 2013.  161p.  ISBN: 9788535271621. }

\def\PUDautor{Geraldo Luis Bezerra Ramalho}

\def\PUDdata{2013-04-17}

\def\PUDrevisao{3}

\def\PUDrevisor{José Ciro dos Santos}

\def\PUDdatarevisao{2019-05-15}

\PUDcorpo
